\documentclass[11pt]{article}
\usepackage[a4paper,margin=1in]{geometry}
\usepackage{amsmath,amssymb,amsthm,booktabs,enumitem,mathtools}
\usepackage[T1]{fontenc}
\usepackage{lmodern}
\usepackage{hyperref}

\hypersetup{colorlinks=true,linkcolor=blue,citecolor=blue,urlcolor=blue}
\newtheorem{theorem}{Theorem}[section]
\newtheorem{proposition}[theorem]{Proposition}
\newtheorem{lemma}[theorem]{Lemma}
\newtheorem{corollary}[theorem]{Corollary}
\theoremstyle{definition}
\newtheorem{definition}[theorem]{Definition}
\theoremstyle{remark}
\newtheorem{remark}[theorem]{Remark}
\newcommand{\Mob}{\mu}
\newcommand{\Z}{\mathbb Z}
\newcommand{\dd}{\delta}
\newcommand{\eps}{\varepsilon}

\title{The all-clock supremum for universally safe\
one-site M\"obius capacity factors}
\author{RH research program}
\date{August 2026}

\begin{document}
\maketitle

\begin{abstract}
For every fixed finite clock $q$, we optimize all phase/current-input sign
factors that satisfy the RH-366 distance-two constraint for every possible
ternary input word.  Their limiting M\"obius correlation exists, and its
maximum absolute value is the squarefree-density weighted independent-set
value
\[
 F(q)=\max_{I\cap(I+2)=\varnothing}\sum_{r\in I}\dd_{q,r}.
\]
The clock need not be a minimal period.  Consequently $q\mid Q$ implies
$F(q)\le F(Q)$.  We then prove a special saturation law for the RH-374 square
clocks $q_y=4\prod_{i\le y}p_i^2$: if $q_y\mid Q$ and $Q$ has the same prime
support, then $F(Q)=F(q_y)=B_y$.  The proof uses zero phases forced by $4$
and $9$ to split the two parity cycles; it is not a general cyclic-cover MWIS
law.  Lifting an arbitrary finite clock to such a same-support multiple gives
\[
             \sup_{q<\infty}F(q)=B_\infty,
             \qquad F(q)<B_\infty\quad(q<\infty).
\]
Thus the RH-374 Euler-product floor is exactly the nonattained all-clock
supremum in this one-site class.  Memory-dependent observables, growing
clocks, infinite selectors, adaptive-capacity convergence, intrinsic
operators, traces, Riemann-zero identification, and RH remain outside scope.
\end{abstract}

\section{Frozen capacity and factor class}

The RH-366 distance-two capacity is
\begin{equation}
 K_N=\max_{\substack{\eps_1,\ldots,\eps_N\in\{-1,+1\}\\
 \text{not }(\eps_n,\eps_{n+2})=(+1,+1)}}
 \left|\sum_{n\le N}\Mob(n)\eps_n\right|.
 \label{eq:capacity}
\end{equation}
Its optimizer may read the entire finite prefix and may change with $N$.
RH-366 proves the bracket
\begin{equation}
 \frac4{\pi^2}\le\liminf_{N\to\infty}\frac{K_N}{N}
 \le\limsup_{N\to\infty}\frac{K_N}{N}\le\frac6{\pi^2},
 \label{eq:capacity-bracket}
\end{equation}
but not convergence \cite{RH366}.  RH-371 gives an exact eight-run formula
for every prefix while leaving the required M\"obius run envelope open
\cite{RH371}.  RH-373 and RH-374 instead construct prescribed one-site
certificates, culminating in a strictly increasing square-clock family
$B_y\uparrow B_\infty$ \cite{RH373,RH374}.

\begin{definition}[One-site factor and universal safety]
Fix $q\ge1$.  A $q$-periodic one-site phase/current-input factor is a family
\[
       g_r:\{-1,0,+1\}\longrightarrow\{-1,+1\},
       \qquad r\in\Z/q\Z.
\]
The declared period need not be minimal.  For a ternary word $(a_n)$, its
output is $g_{n\bmod q}(a_n)$.  The factor is \emph{universally safe} if this
output never has two $+1$ values at distance two, for every ternary input
word.  Its active set is
\[
 S_g=\{r:g_r(a)=+1\text{ for at least one }a\in\{-1,0,+1\}\}.
\]
\end{definition}

Universal safety is equivalent to
\begin{equation}
                       S_g\cap(S_g+2)=\varnothing.
 \label{eq:active-independent}
\end{equation}
Indeed, two active phases at distance two can be activated independently;
the converse is immediate.  This retains loops: for $q=1$ and $q=2$,
addition by two fixes every phase, so only the empty active set is safe.
The class is the one-site boundary isolated in RH-372, not the class of all
finite-memory transducers \cite{RH372}.

\section{Exact optimization at one clock}

For $r\pmod q$, let
\begin{equation}
 \dd_{q,r}=\lim_{N\to\infty}\frac1N
 \#\{n\le N:n\equiv r\pmod q,\ \Mob(n)^2=1\}.
 \label{eq:density}
\end{equation}
The squarefree progression sieve gives this limit.  If
$q=\prod_{p\mid q}p^{a_p}$, its useful exact local form is
\begin{equation}
 \pi^2\dd_{q,r}=
 \begin{cases}
 0,&p^2\mid r\text{ for some }p\text{ with }a_p\ge2,\\[2mm]
 \displaystyle\frac6q
 \prod_{p\mid q}(1-p^{-2})^{-1}
 \prod_{\substack{p\parallel q\\p\mid r}}(1-p^{-1}),&\text{otherwise}.
 \end{cases}
 \label{eq:local-density}
\end{equation}
The second product accounts for a modulus containing only one power of $p$.
It is an exact rational coefficient of $\pi^{-2}$ \cite{Mirsky1948}.

For fixed $q$, Davenport cancellation and finite Fourier inversion give
\[
 \sum_{\substack{n\le N\\n\equiv r\ (q)}}\Mob(n)=o(N).
\]
Splitting the squarefree entries by sign therefore proves, for every
one-site factor,
\begin{equation}
 L_q(g):=\lim_{N\to\infty}\frac1N\sum_{n\le N}
 \Mob(n)g_{n\bmod q}(\Mob(n))
 =\frac12\sum_{r\bmod q}\dd_{q,r}
       \bigl(g_r(+1)-g_r(-1)\bigr).
 \label{eq:factor-limit}
\end{equation}
This is the fixed-clock formula proved in the RH-372 certificate layer
\cite{Davenport1937,RH372}.

\begin{definition}[Fixed-clock optimum]
Let
\[
              F(q)=\max_g |L_q(g)|,
\]
where the maximum is over all universally safe $q$-periodic one-site
factors.
\end{definition}

\begin{theorem}[Weighted phase-MWIS formula]
For every finite $q$,
\begin{equation}
 \boxed{\displaystyle
 F(q)=\max_{\substack{I\subset\Z/q\Z\\I\cap(I+2)=\varnothing}}
                  \sum_{r\in I}\dd_{q,r}.}
 \label{eq:F-MWIS}
\end{equation}
Both correlation orientations attain the value.  In particular,
$F(1)=F(2)=0$.
\end{theorem}

\begin{proof}
Put $c_r=(g_r(+1)-g_r(-1))/2\in\{-1,0,+1\}$.  Whenever $c_r\ne0$, phase
$r$ is active, so the support of $(c_r)$ is contained in the independent set
$S_g$.  Equation \eqref{eq:factor-limit} gives
\[
 |L_q(g)|\le\sum_{r:c_r\ne0}\dd_{q,r}
 \le\max_{I\cap(I+2)=\varnothing}\sum_{r\in I}\dd_{q,r}.
\]
Conversely, for an independent set $I$, set
\[
 g_r(+1)=+1,\qquad g_r(0)=g_r(-1)=-1\quad(r\in I),
 \qquad g_r\equiv-1\quad(r\notin I).
\]
This universally safe factor has $L_q(g)=\sum_{r\in I}\dd_{q,r}$.  Reversing
the roles of $+1$ and $-1$ on $I$ realizes the negative orientation.  The
self-loops at $q=1,2$ force $I=\varnothing$.
\end{proof}

For odd $q>1$, addition by two is one $q$-cycle; for even $q>2$, it is two
cycles of length $q/2$.  Thus \eqref{eq:F-MWIS} includes odd clocks rather
than silently reducing to the two-parity case used by the square family.

\begin{remark}[Deletable phases]
The value $g_r(0)$ contributes nothing to \eqref{eq:factor-limit}; changing
an unnecessary $+1$ there to $-1$ can only remove a safety constraint.
Likewise, an active zero-density phase can be made inactive without changing
the objective.  These observations do not delete positive-weight phases and
do not enlarge the theorem beyond one-site factors.
\end{remark}

\section{Divisibility monotonicity}

\begin{lemma}[Density aggregation]
If $q\mid Q$, then for every $r\pmod q$,
\begin{equation}
              \dd_{q,r}=\sum_{\substack{s\pmod Q\\s\equiv r\ (q)}}
                         \dd_{Q,s}.
 \label{eq:aggregation}
\end{equation}
\end{lemma}

\begin{proof}
For every finite $N$, the squarefree integers in the class $r\pmod q$ are
the disjoint union of the classes $s\pmod Q$ above it.  Divide by $N$ and
take the finitely many limits.
\end{proof}

\begin{proposition}[Clock divisibility]
\label{prop:clock-divisibility}
If $q\mid Q$, then
\begin{equation}
                              F(q)\le F(Q).
 \label{eq:monotonicity}
\end{equation}
\end{proposition}

\begin{proof}
View a $q$-periodic factor as the $Q$-periodic factor
$\widetilde g_s=g_{s\bmod q}$.  This is why minimal period was not required.
Condition \eqref{eq:active-independent} shows that universal safety is
preserved.  The two factors produce the same output word, hence the same
limit; equivalently, this follows by inserting \eqref{eq:aggregation} into
\eqref{eq:factor-limit}.  Maximizing gives \eqref{eq:monotonicity}.
\end{proof}

\section{Special square-clock saturation}

Let $3=p_1<p_2<\cdots$ be the odd primes and recall the RH-374 quantities
\begin{equation}
 P_y=\prod_{i\le y}p_i^2,\qquad q_y=4P_y,\qquad
 A_y=\prod_{i\le y}(p_i^2-1).
 \label{eq:square-clock}
\end{equation}
On the odd phase cycle, let $O_y$ be the number of odd-length positive runs
in one period of the word
\[
            w_y(k)=\mathbf1_{\{p_i^2\nmid2k+1\ \text{for all }i\le y\}}.
\]
RH-374 proves
\begin{equation}
             F(q_y)=B_y=\frac{4+2O_y/A_y}{\pi^2},
 \label{eq:B-y}
\end{equation}
and proves $B_y$ strictly increasing.  With
\[
 e_m=\prod_{p\ \mathrm{odd}}(1-m/p^2),\qquad e_9=0,
\]
its limit is
\begin{equation}
 B_\infty=\frac{4+2C}{\pi^2},\qquad
 C=\frac1{e_1}\sum_{j\in\{1,3,5,7\}}
       (e_j-2e_{j+1}+e_{j+2}).
 \label{eq:B-infinity}
\end{equation}

\begin{proposition}[Same-prime-support saturation]
\label{prop:saturation}
Suppose $q_y\mid Q$ and $Q$ has exactly the same prime divisors as $q_y$.
Write $Q=Rq_y$.  Then
\begin{equation}
                         \boxed{F(Q)=F(q_y)=B_y.}
 \label{eq:saturation}
\end{equation}
\end{proposition}

\begin{proof}
Because every prime in the support occurs to exponent at least two, a phase
$s\pmod Q$ has positive squarefree density precisely when no supported
prime square divides $s$.  All positive phases have the same weight.  By
\eqref{eq:local-density}, it is $1/R$ times the positive $q_y$ weight, hence
\begin{equation}
                         \dd_{Q,s}=\frac{2}{RA_y\pi^2}.
 \label{eq:Q-unit}
\end{equation}

It remains to count the largest independent subset of the positive support.
The distance-two graph has one even and one odd phase cycle.  On the even
cycle, every phase divisible by $4$ has zero weight; these zero phases split
the support, and all $RA_y$ positive even phases may be selected.  On the
odd cycle, the support is $R$ repetitions of the $q_y$ support.  The factor
$p_1^2=9$ inserts a zero every nine successive odd-cycle sites, so all
positive runs have length at most eight and the repetitions do not join
across a hidden all-positive seam.  The RH-374 run count therefore repeats
literally: its odd-cycle MWIS cardinality is $R(A_y+O_y)$.  The total support
MWIS cardinality is
\begin{equation}
                         R(2A_y+O_y).
 \label{eq:scaled-count}
\end{equation}
Multiplying \eqref{eq:scaled-count} by \eqref{eq:Q-unit} gives
$(4+2O_y/A_y)/\pi^2=B_y$.
\end{proof}

\begin{remark}[No general cover law]
An MWIS need not scale under an arbitrary cyclic cover.
Proposition~\ref{prop:saturation}
uses the particular squarefree support: forced $4$-zero phases split the
even component and forced $9$-zero phases split the odd component.  No claim
is made for a general weighted cyclic graph.
\end{remark}

\section{The all-clock supremum}
\label{sec:all-clock}

\begin{theorem}[Exact finite-clock one-site supremum]
Over all finite clocks and universally safe one-site phase/current-input
factors,
\begin{equation}
                 \boxed{\displaystyle\sup_{q<\infty}F(q)=B_\infty.}
 \label{eq:all-clock}
\end{equation}
The supremum is not attained: every fixed finite $q$ satisfies
\begin{equation}
                              F(q)<B_\infty.
 \label{eq:nonattainment}
\end{equation}
\end{theorem}

\begin{proof}
Fix an arbitrary finite $q$.  Choose $y$ so that every odd prime divisor of
$q$ occurs among $p_1,\ldots,p_y$; when $q$ has no odd prime divisor, take
$y=1$.  Put
\[
                         Q=\operatorname{lcm}(q,q_y).
\]
Then $q\mid Q$, $q_y\mid Q$, and $Q$ has exactly the prime support of $q_y$:
extra exponents in $q$, including high powers of $2$ or of an odd prime,
only enlarge the multiplier $R=Q/q_y$.
Propositions~\ref{prop:clock-divisibility} and~\ref{prop:saturation} give
\[
                         F(q)\le F(Q)=B_y<B_\infty,
\]
where the last inequality uses the strict RH-374 sequence.  This proves the
upper bound in \eqref{eq:all-clock} and the nonattainment statement.
Conversely $F(q_y)=B_y$ for every $y$, so
$\sup_qF(q)\ge\sup_yB_y=B_\infty$.
\end{proof}

\begin{corollary}[Relation to adaptive capacity]
\begin{equation}
                  B_\infty\le
                  \liminf_{N\to\infty}\frac{K_N}{N}.
 \label{eq:capacity-floor}
\end{equation}
\end{corollary}

\begin{proof}
Every universally safe factor supplies an admissible word in
\eqref{eq:capacity}, so $F(q)\le\liminf K_N/N$ for each fixed $q$.  Take the
scalar supremum and apply \eqref{eq:all-clock}.
\end{proof}

This recovers the RH-374 floor while adding an exact interpretation: it is
the all-finite-clock supremum in the one-site class.  It does not identify
$B_\infty$ with an ordinary limit of $K_N/N$ and does not choose $q=q(N)$.

\section{Exact executable audit}

The artifact evaluates \eqref{eq:local-density} with exact
\texttt{Fraction} arithmetic and solves each addition-by-two cycle by a
weighted path/cycle DP.  It retains self-loops and handles odd clocks, high
$2$-adic exponents, high odd-prime exponents, zero weights, and both factor
orientations.  Its independent finite checks are:
\begin{itemize}[leftmargin=*]
\item all $8^q$ factor tables for $1\le q\le4$: $4680$ total tables and
  $249$ universally safe tables;
\item all phase subsets for $1\le q\le10$: $2046$ subsets, with every exact
  optimum agreeing with the cycle DP;
\item eight divisibility pairs, checking every density aggregation row, the
  lifted independent set, and $F(q)\le F(Q)$;
\item ten cofinal lifts, including $q=16,27,125,343$ and the largest audited
  lift $343\mid308700$;
\item the RH-374 square rows
  \[
   \pi^2F(36)=4,\qquad \pi^2F(900)=\frac{49}{12},\qquad
   \pi^2F(44100)=\frac{593}{144}.
  \]
\end{itemize}
Representative lift rows are
\begin{center}
\begin{tabular}{@{}rrrrr@{}}
\toprule
$q$ & $Q$ & $R$ & $\pi^2F(q)$ & $\pi^2F(Q)$\\
\midrule
3 & 36 & 1 & $9/4$ & $4$\\
8 & 72 & 2 & $4$ & $4$\\
27 & 108 & 3 & $3$ & $4$\\
125 & 4500 & 5 & $3$ & $49/12$\\
343 & 308700 & 7 & $3$ & $593/144$\\
\bottomrule
\end{tabular}
\end{center}

A bounded reproduction scan over all $1\le q\le256$ has record clocks
$1,3,4,180$ and maximum $\pi^2F(q)=97/24$ at $q=180$.  This finite scan is
not evidence for \eqref{eq:all-clock}; the theorem is the cofinal argument in
Section~\ref{sec:all-clock}.  The result ledger locks twenty-one repository
inputs.  The ARS
integrity audit is repository-locked: it checks claim/source alignment,
artifact provenance, and all seven research failure modes, but makes no fresh
external bibliographic-verification claim.

\section{Route verdict and Gate ledger}

Route A is \texttt{GO}: equations \eqref{eq:F-MWIS},
\eqref{eq:monotonicity}, \eqref{eq:saturation}, and \eqref{eq:all-clock}
form a standalone exact theorem chain.  Route B is
\texttt{STOP\_SCOPED}.  The class excludes memory-dependent observations;
no higher-order M\"obius correlation theorem is supplied.  There is no
$q(N)$, infinite selector, or Davenport estimate uniform in the clock, and
the adaptive capacity limit remains open.

Gates A--E remain false/open.  These factors are externally prescribed
arithmetic certificates, not intrinsic dynamical operators.  This paper does
not construct a Fredholm or spectral determinant, a time-oriented unitary
completion, a self-adjoint generator, a von-Mangoldt prime-power trace, or a
completed-zeta divisor equality.  It does not identify Riemann zeros,
construct a Hilbert--Polya operator, or prove the Riemann Hypothesis.

\section{Conclusion}

The finite-clock one-site question has an exact endpoint: the RH-374
Euler-product constant is its supremum, approached by square clocks and
attained by no finite clock.  The remaining RH-366 problem is genuinely
different.  It requires control of adaptive or memory-dependent arithmetic
structure; another bounded scan or an unproved growing-clock exchange cannot
supply that control.

\bibliographystyle{plain}
\bibliography{references}
\end{document}
